\section{Struktur Program Assembly}
\label{sec:struktur-program}

Program assembly 8086 untuk DOS memiliki kerangka baku. Kerangka ini
memisahkan kode, data, dan \eng{stack} agar assembler dan linker
dapat menetapkan segmen dengan benar. Bagian ini diringkas dari
\cite{sto2001}, BAB V--VI, hlm.~21--26, dan \cite{gunawan2019},
Bagian 4.4, hlm.~32--34.

\subsection{Kerangka Program}

Kerangka program \file{.EXE} yang dipakai pada buku ini ditunjukkan
pada Kode~\ref{lst:kerangka-program}.

\begin{lstlisting}[
    language=MASM,
    caption={Kerangka baku program assembly 8086 untuk DOS.},
    label={lst:kerangka-program}
]
.MODEL SMALL        ; tentukan model memori
.STACK 100h         ; sediakan ruang stack

.DATA               ; ------- segmen data -------
    ; deklarasi variabel dan konstanta

.CODE               ; ------- segmen kode -------
START:              ; label titik masuk
    MOV AX, @DATA   ; inisialisasi DS
    MOV DS, AX

    ; instruksi program

    MOV AH, 4Ch     ; keluar dari program
    MOV AL, 00h
    INT 21h

END START           ; akhiri dan tunjuk titik masuk
\end{lstlisting}

Setiap bagian kerangka tersebut memiliki peran sebagai berikut.

\begin{description}
    \item[\dir{.MODEL}] Menetapkan model memori program. Model
          \code{SMALL} memakai satu segmen kode dan satu segmen data,
          masing-masing maksimum 64 KB.

    \item[\dir{.STACK}] Menyediakan ruang \eng{stack}. Nilai
          \hex{100} berarti 256 byte.

    \item[\dir{.DATA}] Awal segmen data. Variabel dan konstanta
          dideklarasikan di bagian ini.

    \item[\dir{.CODE}] Awal segmen kode. Seluruh instruksi yang
          dieksekusi berada di bagian ini.

    \item[Label] Nama yang menandai suatu lokasi di dalam program.
          Label titik masuk harus sama dengan nama pada direktif
          \dir{END}.

    \item[\dir{END}] Menutup berkas sumber dan menentukan titik awal
          eksekusi.
\end{description}

\subsection{Model Memori}

Direktif \dir{.MODEL} menentukan susunan segmen program. Pilihan model
yang umum ditunjukkan pada Tabel~\ref{tab:model-memori}.

\begin{table}[H]
    \centering
    \caption{Model memori pada direktif \dir{.MODEL}.}
    \label{tab:model-memori}
    \small
    \begin{tabular}{@{}lp{6.4cm}p{4.4cm}@{}}
        \toprule
        \textbf{Model} & \textbf{Keterangan} &
        \textbf{Penggunaan} \\
        \midrule
        \code{TINY} & Kode, data, dan \eng{stack} dalam satu segmen &
        Program \file{.COM} \\
        \addlinespace
        \code{SMALL} & Satu segmen kode dan satu segmen data &
        Program pembelajaran dan program \file{.EXE} kecil \\
        \addlinespace
        \code{MEDIUM} & Kode lebih dari 64 KB, data kurang dari 64 KB &
        Program dengan kode besar \\
        \addlinespace
        \code{COMPACT} & Data lebih dari 64 KB, kode kurang dari 64 KB &
        Program dengan data besar \\
        \addlinespace
        \code{LARGE} & Kode dan data lebih dari 64 KB &
        Program besar \\
        \addlinespace
        \code{HUGE} & Kode, data, dan \eng{array} lebih dari 64 KB &
        Program sangat besar \\
        \bottomrule
    \end{tabular}
\end{table}

Model \code{SMALL} dipakai pada sebagian besar contoh karena cukup
untuk program pembelajaran dan memudahkan inisialisasi segmen data.
Model \code{TINY} dipakai untuk program \file{.COM}.

\subsection{Label dan Komentar}

Label adalah nama yang menandai lokasi instruksi atau data. Aturan
penulisannya adalah sebagai berikut.

\begin{itemize}
    \item Nama label boleh dimulai dengan huruf, \code{@}, \code{\_},
          atau \code{?}.
    \item Nama label tidak boleh mengandung spasi dan tidak boleh
          diawali angka.
    \item Panjang nama label maksimal 31 karakter.
    \item Label instruksi diakhiri tanda titik dua, misalnya
          \code{START:}.
\end{itemize}

Komentar ditulis setelah tanda titik koma (\code{;}). Semua teks
setelah tanda itu diabaikan assembler. Komentar dipakai untuk
menjelaskan maksud instruksi, bukan mengulang bunyi instruksi.

\begin{lstlisting}[language=MASM, numbers=none]
MOV AH, 09h     ; siapkan fungsi cetak string
LEA DX, pesan   ; DX = alamat string
INT 21h         ; panggil layanan DOS
\end{lstlisting}

\subsection{Program COM dan ORG 100h}

Program \file{.COM} memakai model \code{TINY} dan direktif
\dir{ORG 100h}. Direktif ini memberi tahu assembler bahwa program
akan dimuat mulai offset \hex{100}. Ruang 256 byte sebelum program
dipakai DOS untuk \eng{Program Segment Prefix}.

\begin{lstlisting}[
    language=MASM,
    caption={Kerangka program \file{.COM} dengan \dir{ORG 100h}.},
    label={lst:kerangka-com}
]
.MODEL TINY
.CODE
ORG 100h

START:
    ; instruksi program

    INT 20h         ; akhiri program COM

END START
\end{lstlisting}

Pada program \file{.COM}, DOS sudah menetapkan register segmen ke
segmen program. Karena itu, inisialisasi \reg{DS} seperti pada
program \file{.EXE} tidak diperlukan.

\subsection{Inisialisasi Register Segmen Data}

Program \file{.EXE} tidak otomatis memiliki \reg{DS} yang menunjuk
segmen data. Dua instruksi berikut perlu ditulis di awal program.

\begin{lstlisting}[language=MASM, numbers=none]
MOV AX, @DATA       ; alamat segmen data
MOV DS, AX          ; pasang ke DS
\end{lstlisting}

Simbol \code{@DATA} dihasilkan oleh direktif \dir{.MODEL}. Nilainya
adalah nomor segmen tempat \dir{.DATA} berada. Tanpa inisialisasi
ini, instruksi yang mengakses variabel akan membaca alamat yang
salah.

\begin{peringatanbox}
\vspace{1.1em}
Kesalahan yang sering terjadi pada tahap awal:
\begin{enumerate}
    \item lupa menulis \dir{END} atau menulis nama label yang tidak
          sama dengan label titik masuk;
    \item menaruh instruksi di dalam \dir{.DATA} atau data di dalam
          \dir{.CODE};
    \item lupa menginisialisasi \reg{DS} pada program \file{.EXE};
    \item memakai \dir{ORG 100h} pada program \file{.EXE} sehingga
          alamat awal program kacau.
\end{enumerate}
\end{peringatanbox}
